% ============================================================
% فصل ۶ ─ نظرسنجی‌ها و دیدگاه شهروندان عراقی
% ============================================================

% ── خلاصه فصل (اکنون درباره‌ی خودِ این فصل) ────────────────
\begin{chaptersummary}
  بیش از ۴۰ نظرسنجی بزرگ‌مقیاس (\lr{ABC/BBC}، \lr{Gallup}،
  \lr{Arab Barometer}، \lr{Pew}، \lr{IRI} و غیره) طی دو دهه
  با بیش از \ptho{۲۰۰}{۰۰۰} عراقی مصاحبه کرده‌اند.
  پیام تجمیعی آنها: خوش‌بینی اولیه‌ی ۲۰۰۳ به سرخوردگی عمیق
  بدل شد؛ اعتماد به نهادها فروریخت؛ فساد جای امنیت را
  به‌عنوان دغدغه‌ی نخست گرفت ─ اما تقاضا برای دموکراسی
  هرگز به‌طور کامل فروکش نکرد.
\end{chaptersummary}

% ── پل به فصل بعد ─────────────────────────────────────────
\begin{chapterbridge}
  وقتی پروپاگاندا مجوز جنگ را صادر کرد،
  جنگ جنایاتی به بار آورد که حتی حامیانش
  توان دفاع از آن را ندارند.
  فصل بعد به \textbf{ابوغریب، فلوجه، بلک‌واتر}
  و پرسش حقوق بین‌الملل می‌پردازد.
\end{chapterbridge}

\chapter{نظرسنجی‌ها و دیدگاه شهروندان عراقی: تحلیل تفصیلی}

\begin{epigraphbox}
\textit{«وقتی از یک عراقی بپرسید
\textup{«}زندگی‌ات بهتر شده؟\textup{»}
پاسخ بستگی دارد کی و کجا بپرسید.
اما یک چیز ثابت است:
\textbf{هیچ‌کس} از وضع فعلی راضی نیست.»}\\[4pt]
\hfill --- \keyword{گزارش \lr{Gallup} عراق}، ۲۰۲۱
\end{epigraphbox}


% ████████████████████████████████████████████████████████████
%  بخش ۱ ─ مقدمه: بیست سال شنیدن صدای عراقی‌ها
% ████████████████████████████████████████████████████████████
\section{مقدمه: بیست سال شنیدن صدای عراقی‌ها}

بین سال‌های ۲۰۰۳ تا ۲۰۲۳، بیش از ۴۰ نظرسنجی بزرگ‌مقیاس و معتبر
در عراق انجام شده است.
این نظرسنجی‌ها مجموعاً با بیش از \ptho{۲۰۰}{۰۰۰} نفر مصاحبه
کرده‌اند و یکی از غنی‌ترین مجموعه‌داده‌های افکار عمومی در یک
کشور پسا-جنگی را تشکیل می‌دهند.
آنچه در ادامه می‌آید تحلیلی یکپارچه و \textbf{موضوع‌محور}
(نه مؤسسه‌محور) از این داده‌هاست:
ابتدا «لحظه‌ی امید» ۲۰۰۳ را بررسی می‌کنیم،
سپس مسیر فرسایش اعتماد را از امنیت و اقتصاد تا هویت و دموکراسی
دنبال می‌کنیم و در پایان، پنج درس کلان را استخراج خواهیم کرد.

\subsection{نقشه‌ی جامع نظرسنجی‌ها}

\begin{table}[H]
\centering
\caption{فهرست نظرسنجی‌های اصلی انجام‌شده در عراق (۲۰۰۳–۲۰۲۲)}
\label{tab:all-surveys-map}
\small
\begin{tabularx}{\textwidth}{%
  >{\bfseries\scriptsize}p{2.4cm}
  >{\scriptsize\centering\arraybackslash}p{1.6cm}
  >{\scriptsize\centering\arraybackslash}p{0.8cm}
  >{\scriptsize\centering\arraybackslash}p{1.6cm}
  >{\scriptsize\centering\arraybackslash}p{1.8cm}
  >{\scriptsize}X}
\toprule
\textbf{مؤسسه} & \textbf{سال‌ها} & \textbf{موج} &
\textbf{نمونه/موج} & \textbf{پوشش} & \textbf{محورهای اصلی} \\
\midrule
\lr{ABC/BBC/NHK/ARD}   & ۲۰۰۴–۰۹ & ۵  & \ptho{۲}{۰۰۰}–\ptho{۲}{۵۰۰} & ۱۸ استان & حمله، امنیت، اقتصاد، ائتلاف \\
\addlinespace
\lr{Gallup World Poll}  & ۲۰۰۳–۱۲ & ۱۰+ & \ptho{۱}{۰۰۰}+ & ملی & رضایت، امنیت، اشتغال \\
\addlinespace
\lr{Pew Global Attitudes} & ۲۰۰۳–۱۵ & ۸ & \ptho{۱}{۰۰۰}+ & ملی & نگرش به آمریکا، دموکراسی \\
\addlinespace
\lr{IRI}                 & ۲۰۰۳–۱۳ & ۱۲+ & \ptho{۳}{۰۰۰}+ & ملی+استانی & خدمات، عملکرد دولت \\
\addlinespace
\lr{NDI}                 & ۲۰۰۴–۰۹ & ۸  & \ptho{۲}{۰۰۰}+ & ملی & انتخابات، همزیستی \\
\addlinespace
\lr{Arab Barometer}      & ۲۰۰۶–۲۲ & ۷  & \ptho{۱}{۲۰۰}+ & ملی & دموکراسی، اعتماد، ارزش‌ها \\
\addlinespace
\lr{WVS}                 & ۲۰۰۴, ۱۳  & ۲  & \ptho{۲}{۳۰۰}+ & ملی & ارزش‌ها، هویت، مذهب \\
\addlinespace
\lr{Oxford Research Intl.} & ۲۰۰۳–۰۴ & ۳ & \ptho{۲}{۵۰۰}+ & ملی & نخستین پیمایش پس از سقوط \\
\addlinespace
\lr{D3 Systems/KA}       & ۲۰۰۴–۱۱ & ۲۰+ & \ptho{۲}{۰۰۰}+ & ملی & برای وزارت خارجه و ارتش آمریکا \\
\addlinespace
\lr{ICRSS Baghdad}       & ۲۰۰۸–۲۲ & ۱۰+ & \ptho{۱}{۵۰۰}+ & بغداد & حکمرانی شهری \\
\addlinespace
\lr{Zogby International}  & ۲۰۰۴–۰۸ & ۴  & \ptho{۱}{۰۰۰}+ & ملی & ادراک فرهنگی \\
\addlinespace
\lr{UNDP Iraq}            & ۲۰۰۴, ۱۱ & ۲  & \ptho{۲۰}{۰۰۰}+ & ملی & فقر، اشتغال، خدمات \\
\bottomrule
\end{tabularx}
\end{table}

\didyouknow{%
  مجموع حجم نمونه‌ی تمام نظرسنجی‌های عراق بیش از
  \textbf{\ptho{۲۰۰}{۰۰۰} مصاحبه‌ی حضوری} است ─
  در کشوری که بخش‌هایی از آن سال‌ها میدان جنگ بود.
  پرسشگران محلی اغلب با خطر جانی مواجه بودند:
  دست‌کم دو پرسشگر \lr{D3 Systems} هنگام کار میدانی کشته شدند.
}

\subsection{هشت یافته‌ی کلان}

پیش از ورود به جزئیات، هشت یافته‌ی تکرارشونده در
\textbf{همه‌ی} نظرسنجی‌ها قابل شناسایی است:

\begin{enumerate}[rightmargin=1cm, itemsep=6pt]

\item \textbf{شکاف فرقه‌ای–قومی عمیق و پایدار:}
هویت فرقه‌ای–قومی (شیعه/سنی/کرد) بزرگ‌ترین متغیر تعیین‌کننده‌ی
نگرش‌هاست. تفاوت پاسخ‌ها بین سنی‌ها و کردها گاهی تا ۷۰ درصد است.

\item \textbf{سیر نزولی خوش‌بینی:}
خوش‌بینی اولیه (۲۰۰۳–۰۵) → سرخوردگی (۲۰۰۶–۰۸) → بهبود جزئی
(۲۰۰۹–۱۲) → فروپاشی مجدد با داعش → اوج سرخوردگی در تشرین ۲۰۱۹.

\item \textbf{بی‌اعتمادی فزاینده به نهادها:}
اعتماد به دولت و پارلمان از حدود ۴۵–۵۰٪ (۲۰۰۴) به کمتر از ۱۰–۱۵٪
(۲۰۱۹) کاهش یافت.

\item \textbf{تقاضای پایدار برای دموکراسی:}
حمایت از اصل دموکراسی هرگز به زیر ۴۸٪ نرفت ─ حتی در بدترین سال‌ها.

\item \textbf{فساد به‌مثابه دغدغه‌ی اول:}
از ≈۲۰۱۰ به بعد، فساد جای امنیت را در صدر نگرانی‌ها گرفت.

\item \textbf{مخالفت فزاینده با حضور خارجی:}
تقاضا برای خروج نیروها از ۴۰٪ (۲۰۰۴) به بیش از ۷۰٪ (۲۰۰۸) رسید.

\item \textbf{نارضایتی اقتصادی مزمن:}
بیش از ۶۰٪ عراقی‌ها وضعیت اقتصادی را «بد» یا «بسیار بد» ارزیابی کرده‌اند.

\item \textbf{شکاف نسلی فزاینده:}
جوانان (۱۸–۳۰) ناامیدتر، انتقادی‌تر و خواستار تغییر بنیادین‌اند.

\end{enumerate}


% ████████████████████████████████████████████████████████████
%  بخش ۲ ─ لحظه‌ی امید: نخستین نظرسنجی‌ها (۲۰۰۳)
% ████████████████████████████████████████████████████████████
\section{لحظه‌ی امید: نخستین نظرسنجی‌ها پس از سقوط (۲۰۰۳)}

\noindent
تنها شش ماه پس از سقوط صدام، \lr{Oxford Research International}
نخستین نظرسنجی بزرگ‌مقیاس را در عراق انجام داد.
این لحظه‌ی منحصربه‌فرد ─ پیش از بحران‌های بعدی ─ تصویری
از خوش‌بینی محتاطانه ارائه می‌دهد که مقایسه‌اش
با داده‌های شانزده سال بعد، تکان‌دهنده است.

\begin{surveycard}{\lr{Oxford Research International (ORI)} ─ ۲۰۰۳–۲۰۰۴}
\begin{tabularx}{\textwidth}{>{\bfseries\small}r >{\small}X}
موج‌ها & اکتبر ۲۰۰۳ · فوریه ۲۰۰۴ · نوامبر ۲۰۰۴ \\
حجم نمونه & \ptho{۲}{۵۰۰}+ در هر موج \\
سفارش‌دهنده & \lr{BBC} · \lr{ABC News} · \lr{NHK} · \lr{ARD} \\
ویژگی خاص & \textbf{نخستین} نظرسنجی‌های بزرگ‌مقیاس پس از سقوط رژیم \\
\end{tabularx}
\end{surveycard}

\begin{table}[H]
\centering
\caption{نتایج کلیدی نخستین نظرسنجی پس از سقوط (\lr{ORI} ─ اکتبر ۲۰۰۳)}
\label{tab:ori-first}
\small
\begin{tabularx}{\textwidth}{>{\bfseries\small}r >{\small}X >{\centering\arraybackslash\small}p{2cm}}
\toprule
\textbf{پرسش} & \textbf{متن اصلی} & \textbf{پاسخ مثبت} \\
\midrule
دموکراسی می‌خواهم & \lr{I want democracy for Iraq} & ۷۲٪ \\
خوش‌بین به آینده & \lr{Things will be better in a year} & ۶۷٪ \\
از سرنگونی صدام خوشحالم & \lr{I'm glad Saddam is gone} & ۶۲٪ \\
احساس آزادی بیشتر & \lr{I feel freer now} & ۶۰٪ \\
شرایط زندگی بهتر شده & \lr{Life is better now than before the war} & ۵۸٪ \\
اوضاع امنیتی خوب & \lr{Security in my area is good} & ۵۵٪ \\
ائتلاف باید بماند & \lr{Coalition forces should stay until stable} & ۵۰٪ \\
امنیت بزرگ‌ترین مشکل & \lr{Security is the biggest problem} & ۳۵٪ \\
بیکاری بزرگ‌ترین مشکل & \lr{Unemployment is the biggest problem} & ۲۸٪ \\
فساد بزرگ‌ترین مشکل & \lr{Corruption is the biggest problem} & ۸٪ \\
\bottomrule
\end{tabularx}
\end{table}

% ── مق